\section{Introduction}
\label{sec:introduction}

The rapid advancement of computational mechanics and the growing demand for interactive simulation tools have created a fertile ground for innovation at the intersection of numerical methods and digital engineering. Traditional high-fidelity simulations for geometrically nonlinear problems are computationally expensive, often requiring significant time for a single parameter evaluation. This renders them impractical for real-time applications, interactive teaching, and rapid design exploration. The need for a methodology that preserves simulation fidelity while achieving real-time performance motivates the development of a interactive, pedagogically accessible Digital Twin.

\subsection{Research Context}
This project is conducted within the Équipement pour la Performance et le Numérique (EPN04) department at the \acrlong{cnam} in Paris. The overarching objective is to develop a real-time \acrfull{dt} platform that bridges structural mechanics with interactive software, utilizing \acrfull{iga} and projection-based \acrfull{rom} techniques to ensure both geometric accuracy and computational speed.

\subsection{Problematic and Research Gap}
While Reduced Order Modelling (\acrlong{rom}) techniques have been extensively studied for linear structural problems, their application to \textbf{geometrically nonlinear} problems with \textbf{parametric variations} remains an active area of research. The project addresses the following key challenges:
\begin{itemize}
    \item \textbf{Nonlinear Operator Dependence:} The tangent stiffness matrix $\mat{K}_T(\vect{u})$ changes at every Newton--Raphson iteration, requiring repeated projection onto the reduced basis—a computational bottleneck that standard Galerkin projection alone cannot resolve efficiently.
    \item \textbf{Geometric Parameterization:} When geometry acts as a parameter (e.g., varying aspect ratios), the physical domain $\Omega(\vect{\mu})$ changes, invalidating the fixed-mesh assumption of classical ROM approaches.
    \item \textbf{Hyper-reduction for Nonlinear Terms:} Even after projection, evaluating internal forces $\mat{F}_{\text{int}}(\vect{u})$ requires integration over the full mesh. Hyper-reduction techniques (such as DEIM or ECSW) are required to reduce this cost, but their application to IGA-discretized nonlinear problems is not yet well established in pedagogical software.
\end{itemize}

\subsection{Objectives}
The primary objectives of this internship are to:
\begin{enumerate}
    \item \textbf{Develop an IGA-based computational core} capable of solving 2D geometrically nonlinear elasticity problems using NURBS discretization.
    \item \textbf{Construct a projection-based ROM framework} using POD/SVD for basis extraction, extended to handle nonlinear terms via hyper-reduction.
    \item \textbf{Build a real-time web application} serving as both a research validation tool and an interactive platform, enabling users to explore advanced computational mechanics concepts.
\end{enumerate}